\usepackage{xparse}

\usepackage{xltxtra}

\usepackage{etoolbox}

\usepackage{setspace}

\newtoggle{usefont}
\toggletrue{usefont}

% When usefont is false, the default font will be used, which is more compatible with different platforms. 
%\togglefalse{usefont}

% fonts:
\usepackage{fontspec}
\iftoggle{usefont}{
    \setmainfont[Path="V:/Fonts/Eng/",AutoFakeBold=3,AutoFakeSlant=0.25]{BASKVILL.TTF}
    \setsansfont[Path="V:/Fonts/Eng/",AutoFakeBold=3,AutoFakeSlant=0.25]{POORICH.TTF}
    \setmonofont[Path="V:/Fonts/Eng/",AutoFakeBold=3,AutoFakeSlant=0.25]{cour.ttf}
}{}

\usepackage{xeCJK}
\iftoggle{usefont}{
    \setCJKmainfont[Path="V:/Fonts/Zh/",AutoFakeSlant=0.2,AutoFakeBold=3]{STZHONGS.TTF}
    \setCJKsansfont[Path="V:/Fonts/Zh/",AutoFakeSlant=0.2,AutoFakeBold=3]{STKAITI.ttf}
    \setCJKmonofont[Path="V:/Fonts/Zh/",AutoFakeSlant=0.2,AutoFakeBold=3]{simfang.ttf}
}{}

\usepackage{amsmath,amssymb,amsthm,amsfonts,mathrsfs}
\usepackage{mathtools}
\usepackage{physics}
\usepackage{esint}
\usepackage{bm}

\usepackage{xcolor}

\usepackage{colortbl}
\usepackage{makecell}
\usepackage{multirow}
\usepackage{longtable}
\usepackage{booktabs}

\definecolor{backcolor}{rgb}{0.95,0.95,0.92}
\definecolor{codegray}{gray}{0.6}

\definecolor{darkred}{RGB}{139,0,0}
\definecolor{darkpurple}{RGB}{143, 66, 197}

\definecolor{line15}{RGB}{200,179,142}

\definecolor{green1}{RGB}{0, 156, 39}
\definecolor{green2}{RGB}{214, 254, 224}
\definecolor{blue1}{RGB}{163, 172, 204}
\definecolor{blue2}{RGB}{204, 252, 254}

\definecolor{blue3}{HTML}{D1FCFE} % Theorem blue
\definecolor{blue4}{RGB}{72, 141, 246} % Definition blue
\definecolor{blue5}{HTML}{F0FEFF} % Lemma blue
\definecolor{blue6}{RGB}{20, 230, 253} % Criterion blue
\definecolor{blue7}{RGB}{207, 225, 235} % blue7

\usepackage[margin=1.5cm]{geometry}

\setlength{\parindent}{0pt}

\linespread{1.3}

\usepackage{minted}
\usemintedstyle{default}
\setminted{
    %linenos,
    breaklines,
}

\usepackage[most]{tcolorbox}
\usepackage{fontawesome5}

\newtcbox{\highl}[1][red]{on line,
    arc=0pt,outer arc=0pt,colback=#1!10!white,colframe=#1!50!black,
    boxsep=0pt,left=1pt,right=1pt,top=2pt,bottom=2pt,
    boxrule=0pt,bottomrule=1pt,
}

\newtcbox{\chighl}[1][red]{on line,
    arc=7pt,colback=#1!10!white,colframe=#1!50!black,
    before upper={\rule[-3pt]{0pt}{10pt}},boxrule=1pt,
    boxsep=0pt,left=6pt,right=6pt,top=2pt,bottom=2pt
}

% counters:
\newcounter{thmcounter}[subsection]
\setcounter{thmcounter}{0}
\newcounter{dfncounter}[subsection]
\setcounter{dfncounter}{0}
\newcounter{xmpcounter}[subsection]
\setcounter{xmpcounter}{0}
\newcounter{pcdcounter}
\setcounter{pcdcounter}{0}
\newcounter{notecounter}
\setcounter{notecounter}{0}
\renewcommand{\thethmcounter}{\arabic{section}.\arabic{subsection}.\arabic{thmcounter}}
\renewcommand{\thedfncounter}{\arabic{section}.\arabic{subsection}.\arabic{dfncounter}}
\renewcommand{\thexmpcounter}{\arabic{section}.\arabic{subsection}.\arabic{xmpcounter}}

% toggles:
\newtoggle{showremark}
\newtoggle{showproof}
\newtoggle{showvar}
\newtoggle{showxmp}
\toggletrue{showremark}
\toggletrue{showproof}
\toggletrue{showvar}
\toggletrue{showxmp}

\NewDocumentEnvironment{qst}{s O{NaN}}{\begin{tcolorbox}
    [enhanced, colback=white, boxrule=0pt, frame hidden, borderline west={0.7mm}{0.1mm}{darkred}, breakable]
    \textbf{{\Large Question~#2}}\phantomsection\label{qst:#2}
    \IfBooleanTF{#1}{\newline}{}
}{\end{tcolorbox}}

% proof:
\NewDocumentEnvironment{prf}{+b}{
    \iftoggle{showproof}{\begin{tcolorbox}
    [enhanced, colback=lightgray!50!white, boxrule=0pt, frame hidden, breakable, after upper={\hfill$\square$}]
    \textit{Proof:}\newline #1}{}
}{\iftoggle{showproof}{\end{tcolorbox}}{}}

% remark:
\NewDocumentEnvironment{rmk}{O{Contributer} +b}{
    \iftoggle{showremark}{\begin{tcolorbox}
    [enhanced, colback=line15!20!white, boxrule=0pt, frame hidden, fontupper=\sffamily,
        borderline west={0.7mm}{0.1mm}{line15}, breakable, after upper={\hfill\textcolor{codegray}{\textit{remarked~by~#1}}}]
    \textbf{Remark} \newline #2}{}
}{\iftoggle{showremark}{\end{tcolorbox}}{}}

\NewDocumentEnvironment{intrormk}{O{Contributor} +b}{
    \iftoggle{showremark}{\begin{tcolorbox}
    [enhanced, colback=line15!10!white, boxrule=0pt, frame hidden, fontupper=\sffamily,
        borderline north={0.7mm}{0.1mm}{line15}, breakable, after upper={\hfill\textcolor{codegray}{\textit{remarked~by~#1}}}]
    {\textbf{{\centering\large\rmfamily Introduction}}}\newline #2}{}
}{\iftoggle{showremark}{\end{tcolorbox}}{}}

% definition:
\NewDocumentEnvironment{dfn}{s o m s o O{contributor}}{
    \refstepcounter{dfncounter}\begin{tcolorbox}
    [enhanced, colback=white, boxrule=0pt, frame hidden, borderline west={0.7mm}{0.1mm}{blue4}, breakable, after upper={\IfBooleanTF{#4}{\hfill\textcolor{codegray}{\textit{by~#6}}}{}}]
    {\textbf{Definition \thedfncounter :~#3\IfNoValueTF{#5}{}{\IfBlankTF{#5}{}{(#5)}}}}
    \IfBooleanTF{#1}{}{\newline}
}{\end{tcolorbox}\IfNoValueTF{#2}{}{\label{dfn:#2}}}

% theorem:
\NewDocumentEnvironment{thm}{s o m s o O{contributor}}{ % bool(newline), label, name, bool(contributor displayed), subname, contributor
    \refstepcounter{thmcounter}\begin{tcolorbox}
    [enhanced, colback=blue3, boxrule=0pt, frame hidden, borderline west={0.7mm}{0.1mm}{blue1}, breakable, after upper={\IfBooleanTF{#4}{\hfill\textcolor{codegray}{\textit{by~#6}}}{}}]
    {\textbf{Theorem \thethmcounter :~#3\IfNoValueTF{#5}{}{\IfBlankTF{#5}{}{(#5)}}}}
    \IfBooleanTF{#1}{}{\newline}
}{\end{tcolorbox}\IfNoValueTF{#2}{}{\label{thm:#2}}}

\NewDocumentEnvironment{vardfn}{s o m s o O{contributor} +b}{
    \refstepcounter{dfncounter}
    \iftoggle{showvar}{
        \begin{tcolorbox}
        [enhanced, colback=white, boxrule=0pt, frame hidden, borderline west={0.7mm}{0.1mm}{blue4, dotted}, breakable, after upper={\IfBooleanTF{#4}{\hfill\textcolor{codegray}{\textit{by~#6}}}{}}]
        {
            \textbf{{\large Definition \thedfncounter*:~#3 }\IfNoValueTF{#5}{}{\IfBlankTF{#5}{}{(#5)}}}
            \IfBooleanTF{#1}{}{\newline} #7
        }
    }{}
}{\IfNoValueTF{#2}{}{\label{dfn:#2}}\iftoggle{showvar}{\end{tcolorbox}}{}}

\NewDocumentEnvironment{varthm}{s o m s o O{contributor} +b}{
    \refstepcounter{thmcounter}
    \iftoggle{showvar}{\begin{tcolorbox}
    [enhanced, boxrule=0pt, frame hidden, borderline west={0.7mm}{0.1mm}{blue1}, sharp corners, breakable, after upper={\IfBooleanTF{#4}{\hfill\textcolor{codegray}{\textit{by~#6}}}{}}]
    {\textbf{Theorem \thethmcounter*:~#3 \IfNoValueTF{#5}{}{\IfBlankTF{#5}{}{(#5)}}}}
    \IfBooleanTF{#1}{}{\newline} #7}{}
}{\iftoggle{showvar}{\end{tcolorbox}}{}\IfNoValueTF{#2}{}{\label{thm:#2}}}

% property theorem and criterion theorem 
\NewDocumentEnvironment{ppt}{s o m s o O{contributor}}{
    \refstepcounter{thmcounter}\begin{tcolorbox}
    [enhanced, colback=white, boxrule=0pt, frame hidden, borderline west={0.7mm}{0.1mm}{blue1}, breakable, after upper={\IfBooleanTF{#4}{\hfill\textcolor{codegray}{\textit{by~#6}}}{}}]
    {\textbf{{Property Theorem} \thethmcounter~#3\IfNoValueTF{#5}{}{\IfBlankTF{#5}{}{(#5)}}:}}
    \IfBooleanTF{#1}{\newline}{}
}{\end{tcolorbox}\IfNoValueTF{#2}{}{\label{ppt:#2}}}

\NewDocumentEnvironment{crt}{s o m s o O{contributor}}{
    \refstepcounter{thmcounter}\begin{tcolorbox}
    [enhanced, colback=white, boxrule=0pt, frame hidden, 
    borderline west={0.7mm}{0.1mm}{blue6}, interior style={left color=white, right color=blue6!50!white}, 
    breakable, after upper={\IfBooleanTF{#4}{\hfill\textcolor{codegray}{\textit{by~#6}}}{}}]
    {\textbf{{Criterion Theorem} \thethmcounter~#3\IfNoValueTF{#5}{}{\IfBlankTF{#5}{}{(#5)}}:}}
    \IfBooleanTF{#1}{\newline}{}
}{\end{tcolorbox}\IfNoValueTF{#2}{}{\label{crt:#2}}}

% corollary:
\NewDocumentEnvironment{crl}{s o m s o O{contributor}}{
    \refstepcounter{thmcounter}\begin{tcolorbox}
    [enhanced, colback=white!95!blue, boxrule=1pt, frame hidden, 
    borderline west={0.7mm}{0.1mm}{blue1}, 
    breakable, after upper={\IfBooleanTF{#4}{\hfill\textcolor{codegray}{\textit{by~#6}}}{}}]
    {\textbf{{Corollary} \thethmcounter~#3\IfNoValueTF{#5}{}{\IfBlankTF{#5}{}{(#5)}}:}}
    \IfBooleanTF{#1}{}{\newline}
}{\end{tcolorbox}\IfNoValueTF{#2}{}{\label{crl:#2}}}

% exmaple and counterexample:
\NewDocumentEnvironment{xmp}{s o m s o O{contributor} +b}{
    \refstepcounter{xmpcounter}
    \iftoggle{showxmp}{\begin{tcolorbox}
    [enhanced, colback=white, boxrule=0pt, frame hidden, borderline west={0.7mm}{0.1mm}{green1}, breakable, after upper={\IfBooleanTF{#4}{\hfill\textcolor{codegray}{\textit{by~#6}}}{}}]
    {\textbf{{Example} \thexmpcounter~#3\IfNoValueTF{#5}{}{\IfBlankTF{#5}{}{(#5)}}:}}
    \IfBooleanTF{#1}{}{\newline}#7}{}
}{\iftoggle{showxmp}{\end{tcolorbox}}{}\IfNoValueTF{#2}{}{\label{xmp:#2}}}

\NewDocumentEnvironment{cxmp}{s o m s o O{contributor} +b}{
    \refstepcounter{xmpcounter}
    \iftoggle{showxmp}{\begin{tcolorbox}
    [enhanced, colback=green1!25!white, boxrule=0pt, frame hidden, 
    borderline west={0.7mm}{0.1mm}{green1}, breakable, after upper={\IfBooleanTF{#4}{\hfill\textcolor{codegray}{\textit{by~#6}}}{}}]
    {\textbf{{Counterexample} \thexmpcounter~#3\IfNoValueTF{#5}{}{\IfBlankTF{#5}{}{(#5)}}:}}
    \IfBooleanTF{#1}{}{\newline}#7}{}
}{\iftoggle{showxmp}{\end{tcolorbox}}{}\IfNoValueTF{#2}{}{\label{cxmp:#2}}}

% instance:
\NewDocumentEnvironment{ins}{s o m s o O{contributor} +b}{
    \refstepcounter{xmpcounter}
    \iftoggle{showxmp}{\begin{tcolorbox}
    [enhanced, colback=green2!50!white, boxrule=0pt, frame hidden, borderline west={0.7mm}{0.1mm}{green1}, breakable, after upper={\IfBooleanTF{#4}{\hfill\textcolor{codegray}{\textit{by~#6}}}{}}]
    {\textbf{{Instance} \thexmpcounter~~#3\IfNoValueTF{#5}{}{\IfBlankTF{#5}{}{(#5)}}:}}
    \IfBooleanTF{#1}{\newline}{} #7}{}
}{\iftoggle{showxmp}{\end{tcolorbox}}{}\IfNoValueTF{#2}{}{\label{ins:#2}}}

% ========================Reused Envs==========================

\NewDocumentEnvironment{note}{o m o}{\begin{tcolorbox}
    [enhanced, coltext=black, colback=white!95!yellow, colframe=black, boxrule=0.4mm, arc=1mm, 
    before upper={\hspace{2mm}\faQuoteLeft~}, after upper={~\faQuoteRight}, breakable, 
    title=Note \thenotecounter : #2\IfNoValueTF{#1}{}{\phantomsection\label{note:#1}}, fonttitle=\bfseries\large, left=1em, right=1em, top=0.5em, bottom=0.5em]
    {\stepcounter{notecounter}\IfNoValueTF{#3}{}{\IfBlankTF{#3}{}{\textbf{#3}:\newline}}}
}{\end{tcolorbox}}

\NewDocumentEnvironment{varnote}{o m o}{\begin{tcolorbox}
    [enhanced, coltext=black, colback=white!95!blue, colframe=blue1, boxrule=0.4mm, arc=1mm, 
    before upper={\hspace{2mm}\faQuoteLeft~}, after upper={~\faQuoteRight}, breakable, 
    title=Note \thenotecounter : #2\IfNoValueTF{#1}{}{\phantomsection\label{note:#1}}, fonttitle=\bfseries\large, left=1em, right=1em, top=0.5em, bottom=0.5em]
    {\stepcounter{notecounter}\IfNoValueTF{#3}{}{\IfBlankTF{#3}{}{\textbf{#3}\newline}}}
}{\end{tcolorbox}}

\usepackage{graphicx}

\usepackage{subcaption}

\usepackage{enumitem}

\usepackage{hyperref}

\hypersetup{
    colorlinks=true,
    linkcolor=blue,
    filecolor=magenta,
    urlcolor=cyan
}

% Tikz
\usepackage{tikz}
\usetikzlibrary{shapes,arrows,positioning,calc,decorations.pathreplacing,patterns,3d,intersections}

\DeclareMathOperator{\N}{\mathbb{N}}                    % Set of Natural Numbers
\DeclareMathOperator{\Z}{\mathbb{Z}}                    % Set of Integers
\DeclareMathOperator{\Q}{\mathbb{Q}}                    % Set of Rational Numbers
\DeclareMathOperator{\R}{\mathbb{R}}                    % Set of Real Numbers
\DeclareMathOperator{\C}{\mathbb{C}}                    % Set of Complex Numbers
\DeclareMathOperator{\0}{\mathbf{0}}
\DeclareMathOperator{\x}{\mathbf{x}}
\DeclareMathOperator{\w}{\mathbf{w}}
%\DeclareMathOperator{\ii}{\mathrm{i}}
\DeclareMathOperator{\ee}{\mathrm{e}}
\DeclareMathOperator{\st}{\text{s.t.}}
\DeclareMathOperator{\B}{\mathcal{B}}                   % Topological Basis
\DeclareMathOperator{\Log}{Log}

\DeclareMathOperator{\card}{card}
\DeclareMathOperator{\dist}{dist}
\DeclareMathOperator{\spn}{span}
\DeclareMathOperator{\interior}{int}                     % Interior
\DeclareMathOperator{\cl}{cl}                            % Closure

\DeclareMathOperator{\Id}{Id}                           % Identity
\DeclareMathOperator{\img}{im}                          % Image
\DeclareMathOperator{\coker}{coker}                      % Cokernel
\DeclareMathOperator{\ind}{ind}                      % Index

\DeclareMathOperator{\Hom}{Hom}                         % Set of Homomorphisms
\DeclareMathOperator{\End}{End}                         % Set of Endomorphisms
\DeclareMathOperator{\Aut}{Aut}                         % Set of Automorphisms
\DeclareMathOperator{\Isom}{Isom}                       % Set of Isomorphisms

\DeclareMathOperator{\F}{\mathbb{F}}                    % Number Field (F)
\DeclareMathOperator*{\esssup}{\mathrm{ess\,sup}}              % Essential Supremum

% Complex Analysis Hyperrefs
\newcommand{\ii}{\hyperref[dfn:Complex-Number]{\ensuremath{\mathrm{i}}}}
\newcommand{\barC}{\hyperref[dfn:Extended-Complex]{\ensuremath{\bar{\mathbb{C}}}}}
\newcommand{\Arg}{\hyperref[dfn:Argument]{\ensuremath{\operatorname{Arg}}}}
\newcommand{\Ln}{\hyperref[dfn:Logarithmic]{\ensuremath{\operatorname{Ln}}}}
\newcommand{\Windn}{\hyperref[dfn:Winding-Number]{\ensuremath{n}}}
\renewcommand{\Res}{\hyperref[dfn:Residue]{\ensuremath{\operatorname{Res}}}}

%\newcommand{\复可微}{soo}{\IfNoValueTF{#2}{}{#2\,\ensuremath{\text{在}}}\IfNoValueTF{#3}{}{\,#3\,\ensuremath{\text{处}}}\hyperref[dfn:Complex-Derivative]{\ensuremath{\text{复可微}}}[复可微.]}
\newcommand{\Holo}{\hyperref[dfn:Holomorphic-Function]{\ensuremath{\mathcal{H}}}}
\newcommand{\Mobius}{\hyperref[dfn:Mobius-Transformation]{\mathfrak{M}}}
\newcommand{\D}{\hyperref[dfn:Unit-Disk]{\mathbb{D}}}

\newcommand{\CTranslation}{\hyperref[dfn:Basic-Mobius-Transformation]{\mathcal{T}_{\mathfrak{M}}}}
\newcommand{\CScaling}{\hyperref[dfn:Basic-Mobius-Transformation]{\mathcal{S}_{\mathfrak{M}}}}
\newcommand{\CRotation}{\hyperref[dfn:Basic-Mobius-Transformation]{\mathcal{R}_{\mathfrak{M}}}}
\newcommand{\CInversion}{\hyperref[dfn:Basic-Mobius-Transformation]{\mathcal{I}_{\mathfrak{M}}}}

\NewDocumentCommand{\XRatio}{mmmm}{\left(#1,#2\hyperref[dfn:Cross-Ratio]{;}#3,#4\right)}

\newcommand{\Holomorphic}{\hyperref[dfn:Holomorphic-Function]{全纯函数}}
\newcommand{\Analytic}{\hyperref[dfn:Analyticity]{解析函数}}
\newcommand{\Meromorphic}{\hyperref[dfn:Meromorphic-Function]{亚纯函数}}

\newcommand{\AngularForm}{\hyperref[dfn:Angular-Form]{三角表示}}
\newcommand{\PowerSeries}{\hyperref[dfn:Power-Series]{幂级数}}
\newcommand{\UniformConvergent}{\hyperref[dfn:Uniform-Convergence]{一致收敛}}
\newcommand{\ConformalMapping}{\hyperref[dfn:Conformal-Mapping]{共形映照}}
\newcommand{\EssentialSingularity}{\hyperref[dfn:Essential-Singularity]{本性奇点}}
\newcommand{\Taylor}{\hyperref[thm:Taylor]{Taylor 定理}}
\newcommand{\CauchyThm}{\hyperref[thm:Cauchy]{Cauchy积分定理}}

% Colored:
\newcommand{\colora}[1]{\textcolor{red}{#1}}
\newcommand{\colorb}[1]{\textcolor{blue4}{#1}}
\newcommand{\colorc}[1]{\textcolor{blue6}{#1}}
\newcommand{\colord}[1]{\textcolor{darkpurple}{#1}}

\newcommand{\continued}{{\Large\textbf{To be Continued...}}}
\newcommand{\bysorry}{\colorbox{gray!20!white}{\texttt{\textcolor{black}{:= }\textcolor{blue}{by }\textcolor{red}{\textbf{sorry}}}}\;\faTimesCircle\;\,\texttt{\textcolor{gray}{{\footnotesize declaration uses 'sorry'.}}}}

% signature
\newcommand{\Dedicatia}{{\texttt{\textcolor{darkpurple}{Dedicatia}}}}

\usepackage{fancyhdr}
\pagestyle{fancy}

\fancyhf{}

\renewcommand{\headrulewidth}{0.4pt}
\renewcommand{\footrulewidth}{0.25pt}

\fancyhead[L]{Some Notes about Complex Analysis}
\fancyhead[R]{MATH2609-01}
\fancyfoot[C]{Page.\thepage}

\title{{\huge 复分析} \\ Complex Analysis}
\author{Dedicatiaaaaa}
\date{Jun. 24th, 2026}

% Preamble ends here